% !TEX root = main.tex
%% Figure & Graphics %%%%%%%%%%%%%%%%%%%%%%%%%%%%%%%%%%%%%%%%%%%%%%%%%%%%%%%%%%
\usepackage{graphicx}
\usepackage[dvipsnames]{xcolor}
{% raw -%}
%\graphicspath{{subdir1/}{../dir2}}
{% endraw -%}
% Nord colors
\definecolor{night1}{rgb}{0.180, 0.204, 0.251}
\definecolor{night2}{rgb}{0.231, 0.259, 0.322}
\definecolor{night3}{rgb}{0.263, 0.298, 0.369}
\definecolor{night4}{rgb}{0.298, 0.337, 0.416}
\definecolor{snow1}{rgb}{0.847, 0.871, 0.914}
\definecolor{snow2}{rgb}{0.898, 0.914, 0.941}
\definecolor{snow3}{rgb}{0.925, 0.937, 0.957}
\definecolor{frost1}{rgb}{0.561, 0.737, 0.733}
\definecolor{frost2}{rgb}{0.533, 0.753, 0.816}
\definecolor{frost3}{rgb}{0.506, 0.631, 0.757}
\definecolor{frost4}{rgb}{0.369, 0.506, 0.675}
\definecolor{nordred}{rgb}{0.749, 0.380, 0.416}
\definecolor{nordorange}{rgb}{0.816, 0.529, 0.439}
\definecolor{nordyellow}{rgb}{0.922, 0.796, 0.545}
\definecolor{nordgreen}{rgb}{0.639, 0.745, 0.549}
\definecolor{nordviolet}{rgb}{0.706, 0.557, 0.678}
%
%%%% Math expression %%%%%%%%%%%%%%%%%%%%%%%%%%%%%%%%%%%%%%%%%%%%%%%%%%%%%%%%%%
{% raw -%}
\usepackage{amsmath,amssymb}
\usepackage{mathtools}
\usepackage{bm}
\usepackage{braket}
\usepackage{siunitx}
\sisetup{input-digits = 0123456789\pi}
% \mathtoolsset{showonlyrefs=true} % Number only referenced equations
\DeclareMathOperator{\Tr}{Tr}
\renewcommand{\Re}{\operatorname{Re}}
\renewcommand{\Im}{\operatorname{Im}}
\newcommand{\defas}{\equiv}
% theorem like environment
\usepackage{amsthm}
% --- plainbreak: line-break version of plain (body italic, head bold) ---
\newtheoremstyle{plainbreak}%
  {\topsep}%   space above the head
  {\topsep}%   space below the body
  {\itshape}%  body font
  {}%          indent
  {\bfseries}% head font
  {}%          punctuation after the head
  {\newline}%  line break after the head
  {\thmname{#1}\thmnumber{ #2}\thmnote{: #3}}%
% --- definitionbreak: line-break version of definition (body upright, head bold) ---
\newtheoremstyle{definitionbreak}%
  {\topsep}%
  {\topsep}%
  {\upshape}%  body is upright
  {}%
  {\bfseries}%
  {}%
  {\newline}%
  {\thmname{#1}\thmnumber{ #2}\thmnote{: #3}}%
% --- remarkbreak: line-break version of remark (body upright, head italic, reduced spacing) ---
\newtheoremstyle{remarkbreak}%
  {0.5\topsep}%  remark-type styles use half the vertical spacing
  {0.5\topsep}%
  {\upshape}%
  {}%
  {\itshape}%    head is italic (not bold)
  {}%
  {\newline}%
  {\thmname{#1}\thmnumber{ #2}\thmnote{: #3}}%
% --- lawbreak: for laws (title only as head, no number, line break after the head) ---
\newtheoremstyle{lawbreak}%
  {\topsep}%
  {\topsep}%
  {\upshape}%   body is upright
  {}%
  {\bfseries}%  head bold
  {}%           no punctuation (the title is the head itself)
  {\newline}%
  {\thmnote{#3}}%   <- display only the note as the head
% Law env
\theoremstyle{lawbreak}
\newtheorem*{law@inner}{Law}% internal environment name; "Law" is never shown (since \thmname is unused)
\makeatletter
\NewDocumentEnvironment{law}{ m }{%
  \begin{law@inner}[#1]%
}{%
  \end{law@inner}%
}
\makeatother
% Envs with number
\theoremstyle{plainbreak}
%\theoremstyle{plain} % style without line break
\newtheorem{thm}{Theorem}[section]
\newtheorem{lem}[thm]{Lemma}
\newtheorem{prop}[thm]{Proposition}
\newtheorem{cor}[thm]{Corollary}
\newtheorem{conj}[thm]{Conjecture}
%
\theoremstyle{definitionbreak}
%\theoremstyle{definition} % style without line break
\newtheorem{asm}[thm]{Assumption}
\newtheorem{dfn}[thm]{Definition}
%
\theoremstyle{remark}
\newtheorem{rem}[thm]{Remark}
%
\theoremstyle{definitionbreak}
%\theoremstyle{definition} % style without line break
\newtheorem{exc}{Exercise}[section]
%
% Envs without number
\theoremstyle{plainbreak}
%\theoremstyle{plain} % style without line break
\newtheorem*{thm*}{Theorem}
\newtheorem*{lem*}{Lemma}
\newtheorem*{prop*}{Proposition}
\newtheorem*{cor*}{Corollary}
\newtheorem*{conj*}{Conjecture}
%
\theoremstyle{definitionbreak}
%\theoremstyle{definition} % style without line break
\newtheorem*{asm*}{Assumption}
\newtheorem*{dfn*}{Definition}
%
\theoremstyle{remark}
\newtheorem*{rem*}{Remark}
%
\theoremstyle{definitionbreak}
%\theoremstyle{definition} % style without line break
\newtheorem*{exc*}{Exercise}
%%%% Callout block %%%%%%%%%%%%%%%%%%%%%%%%%%%%%%%%%%%%%%%%%%%%%%%%%%%%%%%%%%%%
\usepackage[most]{tcolorbox}
\tcbset{
  calloutbase/.style={
    arc=0.7mm,
    boxrule=0mm,
    left=2.5mm,
    leftrule=0.7mm,
    right=3.5mm,
    top=1.5mm,
    bottom=1.5mm,
    colback=snow3,
  },
  callout/.style 2 args={   % #1: title, #2: space below the title
    calloutbase,
    coltitle=black,
    fonttitle=\bfseries,
    detach title,
    before upper={\tcbtitle\par\vspace{#2}},
    title={#1},
  },
}
% --- calloutbox: framed box with a title (for direct use) ---
\NewTColorBox{calloutbox}{ O{night4} o O{1mm} }{%
  colframe=#1,
  IfValueTF={#2}{callout={#2}{#3}}{calloutbase}%
}
% --- lookup table: environment name -> frame color ---
\makeatletter
\NewDocumentCommand{\RegisterCalloutEnv}{ m m O{optional} }{%
  % #1: environment name, #2: color, #3: title-passing mode optional / mandatory
  \@namedef{callout@color@#1}{#2}%
  \@namedef{callout@titlemode@#1}{#3}%
}
\makeatother
% Mapping between environments and colors
\RegisterCalloutEnv{thm}{frost1}
\RegisterCalloutEnv{lem}{frost1}
\RegisterCalloutEnv{prop}{night4}
\RegisterCalloutEnv{cor}{night4}
\RegisterCalloutEnv{conj}{night4}
\RegisterCalloutEnv{thm*}{frost1}
\RegisterCalloutEnv{lem*}{frost1}
\RegisterCalloutEnv{prop*}{night4}
\RegisterCalloutEnv{cor*}{night4}
\RegisterCalloutEnv{conj*}{night4}
%
\RegisterCalloutEnv{dfn}{nordviolet}
\RegisterCalloutEnv{asm}{nordviolet}
\RegisterCalloutEnv{dfn*}{nordviolet}
\RegisterCalloutEnv{asm*}{nordviolet}
%
\RegisterCalloutEnv{law}{nordviolet}[mandatory]
%
\RegisterCalloutEnv{rem}{night4}
\RegisterCalloutEnv{rem*}{night4}
%
\RegisterCalloutEnv{exc}{nordgreen}
\RegisterCalloutEnv{exc*}{nordgreen}
%
% --- callout: for wrapping theorem-like environments ---
%     the optional title is passed as the optional argument (theorem note) of the inner environment
\makeatletter
\NewDocumentEnvironment{callout}{ o m }{%
  \ifcsname callout@color@#2\endcsname\else
    \GenericError{}{callout environment: `#2' is not registered.
      Register it with \RegisterCalloutEnv}{}{}%
  \fi
  \begin{calloutbox}[\@nameuse{callout@color@#2}]%
  \ifcsstring{callout@titlemode@#2}{mandatory}
    {% environment that requires a title
     \IfValueTF{#1}
       {\begin{#2}{#1}}
       {\GenericError{}{callout environment: `#2' requires a title}{}{}%
        \begin{#2}{??}}%
    }
    {% environment with an optional title
     \IfValueTF{#1}
       {\begin{#2}[#1]}
       {\begin{#2}}%
    }%
}{%
  \end{#2}%
  \end{calloutbox}%
}
\makeatother
%
\NewDocumentCommand{\NewCalloutBoxEnv}{ m m o }{%
  % #1: environment name, #2: color, #3: default title (optional)
  \IfValueTF{#3}
    {% with a default title: use #3 when the environment's title is omitted
     \NewDocumentEnvironment{#1}{ O{#3} }{%
       \begin{calloutbox}[#2][##1]%
     }{%
       \end{calloutbox}%
     }%
    }
    {% without a default title: a box with no title when omitted
     \NewDocumentEnvironment{#1}{ o }{%
       \IfValueTF{##1}
         {\begin{calloutbox}[#2][##1]}
         {\begin{calloutbox}[#2]}%
     }{%
       \end{calloutbox}%
     }%
    }%
}
%
\NewCalloutBoxEnv{objective}{nordred}[Objective]
\NewCalloutBoxEnv{summary}{nordred}[Summary]
\NewCalloutBoxEnv{info}{nordyellow}
{% endraw -%}
%
%% Bibliography %%%%%%%%%%%%%%%%%%%%%%%%%%%%%%%%%%%%%%%%%%%%%%%%%%%%%%%%%%%%%%%
{% if cookiecutter.bibliography == "biblatex+bibtex" -%}
\usepackage[backend=bibtex, style=phys, biblabel=brackets,
           language=auto, autolang=langname,
           articletitle=false, pageranges=false,
           maxnames=3, maxcitenames=2]{biblatex}
\DefineBibliographyStrings{english}{
   andothers = {\mkbibemph{et\addabbrvspace al\adddot}} % Print et al. in italic
}
%\addbibresource{<bibfile1>.bib} % The extension .bib is required!
%\addbibresource{<bibfile2>.bib}
{% elif cookiecutter.bibliography == "biblatex+biber" -%}
\usepackage[backend=biber, style=phys, biblabel=brackets,
           language=auto, autolang=langname,
           articletitle=false, pageranges=false,
           maxnames=3, maxcitenames=2]{biblatex}
\DefineBibliographyStrings{english}{
   andothers = {\mkbibemph{et\addabbrvspace al\adddot}} % Print et al. in italic
}
% \addbibresource{<bibfile1>.bib} % The extension .bib is required!
% \addbibresource{<bibfile2>.bib}
{% elif cookiecutter.bibliography == "bibtex" -%}
\bibliographystyle{plain} 
% stylename=abbrv,acm,alpha,apalike,ieeetr,plain,siam, or unsrt
{% endif -%}
{% if cookiecutter.bibliography == "bibtex" or cookiecutter.bibliography == "manual" -%}
\usepackage{cite}
{% endif -%}
%
%% Settings for margins %%%%%%%%%%%%%%%%%%%%%%%%%%%%%%%%%%%%%%%%%%%%%%%%%%%%%%%
{% raw -%}
%\usepackage{geometry}
%\geometry{left=30truemm,right=30truemm,textheight=25truecm}
{% endraw -%}
%
%% Misc %%%%%%%%%%%%%%%%%%%%%%%%%%%%%%%%%%%%%%%%%%%%%%%%%%%%%%%%%%%%%%%%%%%%%%%
{% raw -%}
\usepackage{booktabs}
\usepackage[hang,flushmargin,perpage]{footmisc}
% hang=remove indent, flushmargin=remove the space between the mark and the footnote, perpage=footnote numbers reset on each page
\setlength{\footnotemargin}{1em} % adjust the space between the footnote mark and the footnote text
\usepackage{algorithm,algpseudocode}
\usepackage{url}
% \usepackage[%
%     %dvipdfmx,
%     %luatex,
%     hidelinks,%
%     colorlinks=true,%
%     urlcolor=frost4,%
%     linkcolor=nordviolet,%
%     citecolor=nordgreen,%
% ]{hyperref}
{% endraw -%}
%
%% Draft %%%%%%%%%%%%%%%%%%%%%%%%%%%%%%%%%%%%%%%%%%%%%%%%%%%%%%%%%%%%%%%%%%%%%%
{% raw -%}
%\usepackage[color]{showkeys}
%\renewcommand*{\showkeyslabelformat}[1]{%
%\fbox{\parbox{40truemm}{\normalfont\small\ttfamily#1}}}
%\usepackage[]{setspace}
%\doublespacing
%\usepackage{xcolor}
%\newcommand*{\TODO}[1]{{\color{red}(TODO: #1)}}
%\newcommand{\remove}[1]{\textcolor{green}{#1}}
%\newcommand{\add}[1]{\textcolor{Purple}{#1}}
% Replace figures with dummy figures
%\let\originalincludegraphics\includegraphics
%\renewcommand{\includegraphics}[2][width=3.2truein]{\originalincludegraphics[#1]{example-image-plain.pdf}}
{% endraw -%}

